\scanpage{195}\mspage{51}
% archive sheet number: 195

\deleted{\unsure{les}} \inserted{le} foncteur induit $I'_{/i}\longrightarrow I_{/i}$ est $M$-cofinal.

La réciproque n'est pas claire, car \deleted{la condition} \inserted{$u$ est $M$-connexe}, \deleted{dite exprime que} \deleted{\unsure{exprime} seulement que, pour} si et seulement si pour tout $i\in\mathrm{Ob}\,I$ et tout \inserted{$M$-}préfaisceau $F_i$ sur $I_{/i}$ \deleted{provenant de} \uline{induit} \deleted{des \unsure{images} \unsure{essentielles}} \uline{\unsure{par|de} un préfaisceau \inserted{$F$} sur} $I$, l'hom

\[
\underset{I_{/i}}{\underset{\longleftarrow}{\lim}}F_i\longrightarrow\underset{I'_{/i}}{\underset{\longleftarrow}{\lim}}(u_{/i})^*(F_i)
\]

est iso. On ne peut donc affirmer d'emblée que cela implique itou pour \uline{tout} $F_i$ sur $I_{/i}$, que si on sait que $(\unsure{U|u}_{/i})^{\unsure{*|\wedge}}:M(I)\longrightarrow M(I_{/i})$ est ess. surjectif. Il en est ainsi pour tout $i$ quand $I$ est \uline{ordonné}, \deleted{mais} car alors $I_{/i}\longrightarrow I$ est pl. fidèle, donc $M(I_{/i})\longrightarrow M(I)$ pl. fidèle, donc $M(I)\longrightarrow M(I_{/i})$ un foncteur de localisation, \deleted{Par contre, ce n'est pas le cas} en particulier, ess. surjectif.

\deleted{\deleted{\unsure{pour}} En fait, il n'y a pas à supposer forcément (pour la réciproque à (4.8.) que les \inserted{$\alpha_i^*:$} $M(I)\longrightarrow M(I_{/i})$ soient ess. surjectifs, mais (sous réserve d'existence de $\alpha_{i*}$) que \uline{$\alpha_{i*}$ soit conservatif}. Car \deleted{\unsure{pour}} il faut exprimer que \deleted{$\unsure{u_*u^*F}\unsure{\longrightarrow|\longleftarrow}F\unsure{(i)}$} \inserted{$\Gamma(I'_{/i},\unsure{\alpha^*|u^*}F)$} $\longleftarrow\Gamma(I_{/i},F)$ est iso, pour tt $F$ dans $M(I)$, or ceci}
